% TOP_PROBLEM: {"id":30007743,"problem_number":"LOCAL-30007743","title":"Migration and housing affordability","category_id":21,"category":{"id":21,"name":"economics","display_name":"Economics","description":"Open economic theory statements, published research agendas, and proposed scoped specifications; question provenance and historical priority are qualified per record.","slug":"economics","order_index":21,"created_at":"2026-10-08T00:00:00Z"},"status":"research_question","external_url":null,"legacy_ids":[],"published":false,"statement":"Estimate the national welfare effect of one metropolitan density reform after wages, rents, migration, and amenities adjust, and isolate migration’s contribution.","economics":{"original_id":232,"field":"Urban, housing and spatial economics","kind":"empirical","formalization_basis":"adapted_research_question","source_status":"research_agenda","priority_status":"earliest_found","priority_note":"Earliest explicit statement located in this audit; no exhaustive priority search or first-ever claim. The mathematical specification below is adapted, not quoted from the source.","earliest_verified_reference":{"authors":["Nathaniel Baum-Snow","Gilles Duranton"],"title":"Housing supply and housing affordability","year":2025,"url":"https://realestate.wharton.upenn.edu/wp-content/uploads/2025/12/BaumSnow_Duranton_bc_2025-003.pdf","doi":"10.1016/bs.hesreg.2025.05.003","locator":"§7 Conclusions and future research directions, printed pp. 448–449","evidence_summary":"The authors explicitly seek dynamic supply elasticities, construction-productivity causes and remedies, and richer spatial models of migration and housing reform."},"current_reference":{"authors":["Nathaniel Baum-Snow","Gilles Duranton"],"title":"Housing supply and housing affordability","year":2025,"url":"https://realestate.wharton.upenn.edu/wp-content/uploads/2025/12/BaumSnow_Duranton_bc_2025-003.pdf","doi":"10.1016/bs.hesreg.2025.05.003","locator":"§7 Conclusions and future research directions, printed pp. 448–449","evidence_summary":"The authors explicitly seek dynamic supply elasticities, construction-productivity causes and remedies, and richer spatial models of migration and housing reform."},"formalization_note":"The source explicitly identifies the underlying gap or agenda. Population, horizon, interventions, estimands, and auxiliary assumptions are proposed operational choices, not a canonical source theorem. Partial later answers are compatible with this scoped question; the citation alone does not prove absence of every subsequent solution."},"statusQualification":"Published question or research agenda verified at the cited locator. The mathematical specification is adapted for this catalogue and includes additional assumptions; source publication does not establish first-ever priority or certify this exact model remains unsolved. The source explicitly identifies the underlying gap or agenda. Population, horizon, interventions, estimands, and auxiliary assumptions are proposed operational choices, not a canonical source theorem. Partial later answers are compatible with this scoped question; the citation alone does not prove absence of every subsequent solution.","statusReviewedAt":"2026-10-08"}
% FIRST_POSTED: 2026-10-08
% ENTRY_KIND: statement_only
% !TeX program = lualatex
\documentclass[11pt]{article}
\usepackage{fontspec}
\usepackage[a4paper,margin=25mm]{geometry}
\usepackage{amsmath,amssymb,mathtools}
\usepackage{xurl}
\usepackage[hidelinks]{hyperref}
\newcommand{\sref}[2]{\hyperlink{source:#1:#2}{[S#2]}}
\setlength{\emergencystretch}{3em}
\begin{document}
% problemId:30007743
\section*{Migration and housing affordability}\label{30007743}
\subsection{Definitions and mathematical statement}
Consider an interconnected system of metropolitan areas with workers choosing residence and employment. Let \(a\in\{0,1\}\) denote a density reform in one high-productivity area. Each worker \(i\) has baseline location \(r_i\), idiosyncratic location tastes, skills, and moving costs. Under policy \(a\), equilibrium wages \(w_r(a)\), rents \(p_r(a)\), amenities \(A_r(a)\), and residence \(r_i(a)\) are jointly determined. Let \(v_i(a)\) be monetary-equivalent lifetime utility, including moving costs. Define
\[\Delta W=E\!\left[\sum_i\{v_i(1)-v_i(0)\}\right]-K,\]
where \(K\) is reform-associated infrastructure resource cost not already included in utility, and baseline residents, newcomers, and residents elsewhere are all retained. Identify how migration responses mediate rents and welfare, comparing unrestricted equilibrium with a clearly defined counterfactual holding migration choices fixed. Specify labor-demand, housing-supply, and amenity responses, and estimate migration parameters using credible mobility variation. Avoid adding rent transfers to both tenant losses and landlord gains without consistent welfare accounting. The housing review identifies migration assumptions and richer spatial heterogeneity as active research needs. This is a proposed equilibrium estimand adapted from that agenda, not a claim that housing reform always improves local affordability equally.

\par\textbf{Formulation and scope.} The source explicitly identifies the underlying gap or agenda. Population, horizon, interventions, estimands, and auxiliary assumptions are proposed operational choices, not a canonical source theorem. Partial later answers are compatible with this scoped question; the citation alone does not prove absence of every subsequent solution.

\par\textbf{Open-question provenance.} An explicit question or research agenda was verified at the source locator. This does not by itself certify that every later version remains unresolved \sref{30007743}{1}.

\par\textbf{Historical priority.} Earliest explicit statement located in this audit; no exhaustive priority search or first-ever claim. The mathematical specification below is adapted, not quoted from the source.

\subsection{Short English statement}
Estimate the national welfare effect of one metropolitan density reform after wages, rents, migration, and amenities adjust, and isolate migration’s contribution.

\subsection{Sources}
\begin{itemize}\raggedright
\item[\textnormal{[S1]}]\hypertarget{source:30007743:1}{} Nathaniel Baum-Snow, Gilles Duranton. 2025. Housing supply and housing affordability. §7 Conclusions and future research directions, printed pp. 448–449.
\url{https://realestate.wharton.upenn.edu/wp-content/uploads/2025/12/BaumSnow_Duranton_bc_2025-003.pdf}.
DOI: 10.1016/bs.hesreg.2025.05.003.
{\small\textit{Earliest verified question/agenda reference; historical priority qualified below. The authors explicitly seek dynamic supply elasticities, construction-productivity causes and remedies, and richer spatial models of migration and housing reform.}}
\end{itemize}
\end{document}
